% T12 C/D : six sujets d'entraînement de série D, originaux (écrits le 16 novembre 2026 sur le modèle des épreuves du
% baccalauréat série D : deux exercices et un problème, sur 20 points). Les corrigés ne s'impriment que dans l'édition du professeur.
\chapter{Sujets d'entra\^inement, s\'erie D}
\label{chap:sujets entrainement D}
Les six sujets suivants sont construits sur le mod\`ele de l'\'epreuve de math\'ematiques du baccalaur\'eat s\'erie D :
un exercice sur les nombres complexes, un exercice de probabilit\'es ou de statistiques, et un probl\`eme d'analyse.
Dur\'ee conseill\'ee : 4 heures. Le bar\`eme est indiqu\'e pour chaque partie (total : 20 points).
La calculatrice scientifique est autoris\'ee.

%%%%%%%%%%%%%%%%%%%%%%%%%%%%%%%%%%%%%%%%%%%%%%%%%%%%%%%%%%%%%%%%%%%%%%%%%%%%%%%%%%%%%%%%%%%%%%%%%%%%%%%%%
\section{Sujet D1}
\noindent \underline{\textsc{Exercice I}} (5 points) \newline
On consid\`ere le polyn\^ome $P(z) = z^3 - (2+i)z^2 + (2+2i)z - 2i$, o\`u $z \in \mathbb{C}$.
\begin{enumerate}
      \item \begin{enumerate}
            \item Calculer $P(i)$.
            \item D\'eterminer les r\'eels $a$ et $b$ tels que $P(z) = (z - i)(z^2 + az + b)$ pour tout $z$.
            \item R\'esoudre dans $\mathbb{C}$ l'\'equation $P(z) = 0$.
            \end{enumerate}
      \item Dans le plan complexe muni d'un rep\`ere orthonorm\'e direct $(O ; \vec{u}, \vec{v})$ (unit\'e 2~cm),
            on consid\`ere les points $A$, $B$, $C$ d'affixes $z_A = i$, $z_B = 1 + i$, $z_C = 1 - i$.
            Placer les points et calculer $Z = \dfrac{z_C - z_B}{z_A - z_B}$. En d\'eduire la nature du triangle $ABC$.
      \item Soit $S$ la similitude directe de centre $B$, de rapport $2$ et d'angle $\frac{\pi}{2}$.
            \begin{enumerate}
            \item Donner l'\'ecriture complexe de $S$ et v\'erifier que $S(A) = C$.
            \item D\'eterminer l'image par $S$ du cercle $(\Gamma)$ de centre $B$ et de rayon $2$.
            \end{enumerate}
      \item \'Ecrire $z_B$ sous forme trigonom\'etrique, puis calculer $z_B^{2026}$.
\end{enumerate}

\noindent \underline{\textsc{Exercice II}} (5 points) \newline
Un sac contient huit jetons indiscernables au toucher : cinq rouges et trois verts. On tire simultan\'ement trois jetons.
\begin{enumerate}
      \item Calculer la probabilit\'e des \'ev\`enements $A$ : << les trois jetons sont de la m\^eme couleur >> et
            $B$ : << on obtient au moins un jeton vert >>.
      \item Soit $X$ la variable al\'eatoire \'egale au nombre de jetons verts tir\'es.
            Donner la loi de probabilit\'e de $X$, puis calculer son esp\'erance math\'ematique et sa variance.
      \item Un joueur paie 400~Ar pour effectuer un tirage ; il re\c{c}oit 500~Ar par jeton vert tir\'e.
            Soit $G$ son gain alg\'ebrique. Exprimer $G$ en fonction de $X$ et calculer $E(G)$. Le jeu est-il favorable au joueur\,?
            Quelle mise rendrait le jeu \'equitable\,?
      \item Le joueur joue quatre parties ind\'ependantes (les jetons sont remis dans le sac apr\`es chaque partie).
            Calculer la probabilit\'e qu'il tire au moins une fois trois jetons verts (r\'esultat \`a $10^{-3}$ pr\`es).
\end{enumerate}
\noindent\underline{NB :} on donnera les probabilit\'es sous forme de fractions irr\'eductibles, sauf indication contraire.

\noindent \underline{\textsc{Probl\`eme}} (10 points) \newline
On consid\`ere la fonction $f$ d\'efinie sur $\mathbb{R}$ par $f(x) = x - 1 + (x+2)\,\expo{-x}$ et $(\mathscr{C})$ sa courbe
repr\'esentative dans un rep\`ere orthonorm\'e $(O ; \vec{\imath}, \vec{\jmath})$ d'unit\'e 1~cm.

\underline{\textbf{Partie A}} \newline
Soit $g$ la fonction d\'efinie sur $\mathbb{R}$ par $g(x) = 1 - (x+1)\,\expo{-x}$.
\begin{enumerate}
      \item \'Etudier les variations de $g$ (on ne demande pas les limites).
      \item En d\'eduire que $g(x) \geq 0$ pour tout r\'eel $x$, l'\'egalit\'e n'ayant lieu que pour $x = 0$.
\end{enumerate}

\underline{\textbf{Partie B}}
\begin{enumerate}
      \item Calculer les limites de $f$ en $-\infty$ et en $+\infty$.
      \item Montrer que $f'(x) = g(x)$ pour tout r\'eel $x$. Dresser le tableau de variations de $f$.
      \item \begin{enumerate}
            \item Montrer que la droite $(D)$ d'\'equation $y = x - 1$ est asymptote \`a $(\mathscr{C})$ en $+\infty$.
                  \'Etudier la position de $(\mathscr{C})$ par rapport \`a $(D)$.
            \item Calculer $\lim\limits_{x \to -\infty} \frac{f(x)}{x}$ et interpr\'eter.
            \end{enumerate}
      \item Montrer que le point $I(0 ; 1)$ est un point d'inflexion de $(\mathscr{C})$ et donner la tangente en $I$.
      \item Montrer que l'\'equation $f(x) = 0$ admet une unique solution $\alpha$ et que $-1{,}5 < \alpha < -1$.
      \item Tracer $(D)$ et $(\mathscr{C})$.
      \item Soit $\lambda > 0$. \`A l'aide d'une int\'egration par parties, calculer l'aire $\mathscr{A}(\lambda)$, en cm$^2$,
            du domaine limit\'e par $(\mathscr{C})$, $(D)$ et les droites d'\'equations $x = 0$ et $x = \lambda$.
            Calculer $\lim\limits_{\lambda \to +\infty} \mathscr{A}(\lambda)$.
      \item Justifier que $f$ admet une fonction r\'eciproque $f^{-1}$ d\'efinie sur $\mathbb{R}$.
            $f^{-1}$ est-elle d\'erivable en $1$\,? Interpr\'eter graphiquement.
\end{enumerate}
\begin{center}
On donne : $\expo{} \approx 2{,}72$ ; $\expo{1{,}5} \approx 4{,}48$.
\end{center}

%%%%%%%%%%%%%%%%%%%%%%%%%%%%%%%%%%%%%%%%%%%%%%%%%%%%%%%%%%%%%%%%%%%%%%%%%%%%%%%%%%%%%%%%%%%%%%%%%%%%%%%%%
\newpage
\section{Sujet D2}
\noindent \underline{\textsc{Exercice I}} (5 points) \newline
Le plan complexe est muni d'un rep\`ere orthonorm\'e direct $(O ; \vec{u}, \vec{v})$.
\begin{enumerate}
      \item R\'esoudre dans $\mathbb{C}$ l'\'equation $z^2 - 2\sqrt{3}\,z + 4 = 0$.
            On note $z_A$ la solution de partie imaginaire positive et $z_B$ l'autre ; $A$ et $B$ sont les points d'affixes $z_A$ et $z_B$.
      \item \'Ecrire $z_A$ et $z_B$ sous forme exponentielle. Calculer $z_A^6$.
            Montrer que $\frac{z_A}{z_B} = \expo{i\frac{\pi}{3}}$ et en d\'eduire la nature du triangle $OAB$.
      \item Soit $r$ la rotation de centre $O$ et d'angle $\frac{\pi}{3}$. D\'eterminer l'affixe du point $A' = r(A)$.
      \item On consid\`ere l'application $S$ qui \`a tout point $M(z)$ associe $M'(z')$ avec $z' = (1 + i\sqrt{3})\,z$.
            \begin{enumerate}
            \item Pr\'eciser la nature et les \'el\'ements caract\'eristiques de $S$.
            \item Calculer les affixes de $S(A)$ et $S(B)$. Montrer que les points $O$, $A$ et $S(B)$ sont align\'es.
            \end{enumerate}
      \item D\'eterminer et construire l'ensemble $(E)$ des points $M(z)$ tels que $|z - 2i| = |z - \sqrt{3} - i|$.
\end{enumerate}

\noindent \underline{\textsc{Exercice II}} (5 points) \newline
Le tableau suivant donne le nombre d'abonn\'es \`a internet d'une commune, en centaines ;
$x_i$ est le rang de l'ann\'ee ($x = 1$ pour 2020).
\begin{center}
\begin{tabular}{|c|c|c|c|c|c|c|}
\hline
Ann\'ee & 2020 & 2021 & 2022 & 2023 & 2024 & 2025 \\
\hline
Rang $x_i$ & 1 & 2 & 3 & 4 & 5 & 6 \\
\hline
Abonn\'es $y_i$ (centaines) & 52 & 60 & 71 & 79 & 90 & 98 \\
\hline
\end{tabular}
\end{center}
\begin{enumerate}
      \item Repr\'esenter le nuage de points $M_i(x_i ; y_i)$ (1~cm pour une ann\'ee ; 1~cm pour 10 centaines, en commen\c{c}ant \`a 50).
      \item Calculer les coordonn\'ees du point moyen $G$ et le placer.
      \item Calculer la variance de $x$, la covariance de $x$ et $y$ et le coefficient de corr\'elation lin\'eaire $r$
            (\`a $10^{-3}$ pr\`es). Un ajustement affine est-il justifi\'e\,?
      \item D\'eterminer une \'equation de la droite de r\'egression de $y$ en $x$ (coefficients \`a $10^{-2}$ pr\`es) et la tracer.
      \item En supposant que l'\'evolution se poursuit de la m\^eme mani\`ere :
            \begin{enumerate}
            \item estimer le nombre d'abonn\'es en 2027 ;
            \item d\'eterminer \`a partir de quelle ann\'ee le nombre d'abonn\'es d\'epassera 15\,000.
            \end{enumerate}
\end{enumerate}

\noindent \underline{\textsc{Probl\`eme}} (10 points) \newline
\underline{\textbf{Partie A}} \newline
Soit $g$ la fonction d\'efinie sur $]0 ; +\infty[$ par $g(x) = x^2 + \ln x$.
\begin{enumerate}
      \item \'Etudier les variations de $g$ et calculer ses limites aux bornes de son ensemble de d\'efinition.
      \item Montrer que l'\'equation $g(x) = 0$ admet une unique solution $\alpha$ et que $0{,}6 < \alpha < 0{,}7$.
      \item En d\'eduire le signe de $g(x)$ suivant les valeurs de $x$.
\end{enumerate}
\underline{\textbf{Partie B}} \newline
Soit $f$ la fonction d\'efinie sur $]0 ; +\infty[$ par $f(x) = x - \dfrac{1 + \ln x}{x}$, et $(\mathscr{C})$ sa courbe
dans un rep\`ere orthonorm\'e (unit\'e 2~cm).
\begin{enumerate}
      \item Calculer les limites de $f$ en $0$ et en $+\infty$. Interpr\'eter la limite en $0$.
      \item Montrer que $f'(x) = \dfrac{g(x)}{x^2}$ et dresser le tableau de variations de $f$.
      \item Montrer que $f(\alpha) = \dfrac{2\alpha^2 - 1}{\alpha}$.
      \item \begin{enumerate}
            \item Montrer que la droite $(D) : y = x$ est asymptote \`a $(\mathscr{C})$ en $+\infty$.
            \item \'Etudier la position relative de $(\mathscr{C})$ et $(D)$.
            \end{enumerate}
      \item V\'erifier que $f(1) = 0$. Montrer que l'\'equation $f(x) = 0$ admet exactement deux solutions, $1$ et un r\'eel
            $\beta$ tel que $0{,}4 < \beta < 0{,}5$.
      \item Tracer $(D)$ et $(\mathscr{C})$.
      \item Calculer, en cm$^2$, l'aire du domaine limit\'e par $(\mathscr{C})$, $(D)$ et les droites d'\'equations
            $x = \frac{1}{\expo{}}$ et $x = \expo{}$. (On pourra remarquer que $\frac{1 + \ln x}{x}$ est de la forme $u'u$.)
\end{enumerate}
\begin{center}
On donne : $\ln 0{,}4 \approx -0{,}92$ ; $\ln 0{,}5 \approx -0{,}69$ ; $\ln 0{,}6 \approx -0{,}51$ ; $\ln 0{,}7 \approx -0{,}36$ ; $\alpha \approx 0{,}65$.
\end{center}

%%%%%%%%%%%%%%%%%%%%%%%%%%%%%%%%%%%%%%%%%%%%%%%%%%%%%%%%%%%%%%%%%%%%%%%%%%%%%%%%%%%%%%%%%%%%%%%%%%%%%%%%
\newpage
\section{Sujet D3}
\noindent \underline{\textsc{Exercice I}} (5 points) \newline
Dans le plan complexe muni d'un rep\`ere orthonorm\'e direct $(O ; \vec{u}, \vec{v})$ (unit\'e 2~cm), on d\'efinit la suite
de points $M_n$ d'affixes $z_n$ par $z_0 = 4$ et, pour tout entier naturel $n$, $z_{n+1} = \dfrac{1 + i}{2}\, z_n$.
\begin{enumerate}
      \item Calculer $z_1$, $z_2$, $z_3$, $z_4$ et placer les points $M_0$ \`a $M_4$.
      \item \'Ecrire $\frac{1+i}{2}$ sous forme exponentielle. En d\'eduire $z_n$ sous forme exponentielle en fonction de $n$.
      \item On pose $r_n = |z_n|$. Montrer que $(r_n)$ est une suite g\'eom\'etrique et d\'eterminer sa limite.
      \item Calculer $\dfrac{z_n - z_{n+1}}{0 - z_{n+1}}$. En d\'eduire la nature du triangle $O M_{n+1} M_n$.
      \item On pose $\ell_n = M_n M_{n+1}$ et $L_n = \ell_0 + \ell_1 + \cdots + \ell_{n-1}$ (longueur de la ligne bris\'ee $M_0 M_1 \ldots M_n$).
            Exprimer $\ell_n$ puis $L_n$ en fonction de $n$ et calculer $\lim\limits_{n \to +\infty} L_n$.
      \item Pour quelles valeurs de $n$ le point $M_n$ appartient-il \`a la demi-droite $[O ; \vec{u})$\,?
\end{enumerate}

\noindent \underline{\textsc{Exercice II}} (5 points) \newline
Un commer\c{c}ant ach\`ete ses mangues chez trois producteurs : $50\,\%$ chez $P_1$, $30\,\%$ chez $P_2$ et $20\,\%$ chez $P_3$.
Une mangue de $P_1$ est ab\^im\'ee avec la probabilit\'e $0{,}04$, une mangue de $P_2$ avec la probabilit\'e $0{,}05$
et une mangue de $P_3$ avec la probabilit\'e $0{,}1$. On choisit une mangue au hasard dans le stock.
\begin{enumerate}
      \item Construire un arbre pond\'er\'e d\'ecrivant la situation.
      \item Calculer la probabilit\'e que la mangue soit ab\^im\'ee.
      \item La mangue est ab\^im\'ee. Quelle est la probabilit\'e qu'elle provienne de $P_3$\,?
      \item Une cliente ach\`ete trois mangues ; on admet que les trois choix sont ind\'ependants.
            Soit $X$ le nombre de mangues ab\^im\'ees.
            \begin{enumerate}
            \item Quelle est la loi de $X$\,? Calculer $P(X = 0)$ et $P(X \geq 1)$ (\`a $10^{-4}$ pr\`es).
            \item Calculer l'esp\'erance de $X$.
            \end{enumerate}
\end{enumerate}

\noindent \underline{\textsc{Probl\`eme}} (10 points) \newline
Soit $f$ la fonction d\'efinie sur $\mathbb{R}$ par $f(x) = (x+1)^2\,\expo{-x}$ et $(\mathscr{C})$ sa courbe
dans un rep\`ere orthonorm\'e (unit\'e 2~cm).
\begin{enumerate}
      \item Calculer les limites de $f$ en $-\infty$ et en $+\infty$ ; \'etudier la branche infinie en $-\infty$
            et interpr\'eter la limite en $+\infty$.
      \item Montrer que $f'(x) = (x+1)(1-x)\,\expo{-x}$ et dresser le tableau de variations de $f$.
      \item \begin{enumerate}
            \item D\'eterminer une \'equation de la tangente $(T)$ \`a $(\mathscr{C})$ au point d'abscisse $0$.
            \item V\'erifier que $f(x) - (x + 1) = (x+1)(x + 1 - \expo{x})\,\expo{-x}$.
            \item On admet que $\expo{x} \geq x + 1$ pour tout r\'eel $x$, avec \'egalit\'e seulement pour $x = 0$.
                  En d\'eduire la position de $(\mathscr{C})$ par rapport \`a $(T)$.
            \end{enumerate}
      \item Montrer que l'\'equation $f(x) = 1$ admet exactement trois solutions : $0$, un r\'eel $\gamma \in\ ]-2 ; -1[$
            et un r\'eel $\alpha \in\ ]2 ; 3[$.
      \item Tracer $(T)$ et $(\mathscr{C})$.
      \item \begin{enumerate}
            \item V\'erifier que la fonction $F : x \mapsto -(x^2 + 4x + 5)\,\expo{-x}$ est une primitive de $f$ sur $\mathbb{R}$.
            \item Calculer, en cm$^2$, l'aire du domaine limit\'e par $(\mathscr{C})$, l'axe des abscisses et les droites
                  d'\'equations $x = 0$ et $x = 1$.
            \end{enumerate}
      \item Pour tout entier $n \geq 1$, on pose
            \[ u_n = \int_n^{n+1} f(x)\,\mathrm{d}x . \]
            Montrer que $f(n+1) \leq u_n \leq f(n)$ et en d\'eduire la limite de la suite $(u_n)$.
\end{enumerate}
\begin{center}
On donne : $\expo{} \approx 2{,}72$ ; $\expo{2} \approx 7{,}39$ ; $\expo{3} \approx 20{,}09$.
\end{center}

%%%%%%%%%%%%%%%%%%%%%%%%%%%%%%%%%%%%%%%%%%%%%%%%%%%%%%%%%%%%%%%%%%%%%%%%%%%%%%%%%%%%%%%%%%%%%%%%%%%%%%%%
\newpage
\section{Sujet D4}
\noindent \underline{\textsc{Exercice I}} (5 points) \newline
Le plan complexe est muni d'un rep\`ere orthonorm\'e direct $(O ; \vec{u}, \vec{v})$ (unit\'e 1~cm).
\begin{enumerate}
      \item V\'erifier que $(2 + i)^2 = 3 + 4i$, puis r\'esoudre dans $\mathbb{C}$ l'\'equation
            $z^2 - (4 + 3i)\,z + 1 + 5i = 0$. On note $A$ et $B$ les points d'affixes respectives $1 + i$ et $3 + 2i$.
      \item \begin{enumerate}
            \item D\'eterminer l'affixe du point $C$, image de $A$ par la rotation de centre $B$ et d'angle $-\frac{\pi}{2}$,
                  et l'affixe du point $D$, image de $B$ par la rotation de centre $A$ et d'angle $\frac{\pi}{2}$.
            \item Montrer que $ABCD$ est un carr\'e et donner l'affixe de son centre $\Omega$.
            \end{enumerate}
      \item Soit $s$ la similitude directe de centre $A$ qui transforme $B$ en $C$.
            \begin{enumerate}
            \item D\'eterminer l'\'ecriture complexe de $s$, son rapport et son angle.
            \item Calculer les affixes de $s(C)$ et $s(D)$.
            \item Quelle est l'aire de l'image du carr\'e $ABCD$ par $s$\,?
            \end{enumerate}
\end{enumerate}

\noindent \underline{\textsc{Exercice II}} (5 points)
\begin{description}
      \item[I.] Un questionnaire comporte quatre questions ; pour chacune, trois r\'eponses sont propos\'ees dont une seule est juste.
            Un candidat r\'epond au hasard \`a chaque question. Soit $X$ le nombre de bonnes r\'eponses.
            \begin{enumerate}
            \item Justifier que $X$ suit une loi binomiale dont on donnera les param\`etres.
            \item Calculer $P(X \geq 3)$ et la probabilit\'e d'obtenir au moins une bonne r\'eponse.
            \item Une bonne r\'eponse rapporte 3 points, une mauvaise en retire 1. Exprimer la note $N$ en fonction de $X$
                  et calculer son esp\'erance.
            \end{enumerate}
      \item[II.] On a relev\'e, pour six \'el\`eves, le nombre $x$ d'heures de r\'evision et la note $y$ obtenue \`a un devoir.
            \begin{center}
            \begin{tabular}{|c|c|c|c|c|c|c|}
            \hline
            $x_i$ (heures) & 2 & 4 & 5 & 7 & 8 & 10 \\
            \hline
            $y_i$ (note) & 7 & 9 & 10 & 12 & 12 & 15 \\
            \hline
            \end{tabular}
            \end{center}
            \begin{enumerate}
            \item Calculer les coordonn\'ees du point moyen $G$, la covariance de $x$ et $y$ et le coefficient de corr\'elation
                  lin\'eaire (\`a $10^{-2}$ pr\`es).
            \item D\'eterminer la droite de r\'egression de $y$ en $x$ (coefficients \`a $10^{-2}$ pr\`es).
            \item Estimer la note d'un \'el\`eve qui r\'evise 9 heures.
            \end{enumerate}
\end{description}

\noindent \underline{\textsc{Probl\`eme}} (10 points) \newline
Soit $f$ la fonction d\'efinie sur $[0 ; +\infty[$ par $f(x) = x\,(1 - \ln x)^2$ si $x > 0$ et $f(0) = 0$.
On note $(\mathscr{C})$ sa courbe dans un rep\`ere orthonorm\'e (unit\'e 2~cm).
On admet que $\lim\limits_{x \to 0^+} x \ln x = 0$ et $\lim\limits_{x \to 0^+} x (\ln x)^2 = 0$.
\begin{enumerate}
      \item \begin{enumerate}
            \item \'Etudier la continuit\'e de $f$ en $0$.
            \item \'Etudier la d\'erivabilit\'e de $f$ en $0$ et interpr\'eter graphiquement le r\'esultat.
            \end{enumerate}
      \item Calculer $\lim\limits_{x \to +\infty} f(x)$ et $\lim\limits_{x \to +\infty} \frac{f(x)}{x}$. Interpr\'eter.
      \item Montrer que, pour $x > 0$, $f'(x) = (\ln x - 1)(\ln x + 1)$. Dresser le tableau de variations de $f$.
      \item D\'eterminer une \'equation de la tangente $(T)$ \`a $(\mathscr{C})$ au point d'abscisse $1$.
      \item V\'erifier que $f(1) = 1$. Montrer que l'\'equation $f(x) = 1$ admet exactement trois solutions dans $]0 ; +\infty[$.
      \item Tracer $(T)$ et $(\mathscr{C})$.
      \item \begin{enumerate}
            \item \`A l'aide d'une int\'egration par parties, calculer $K = \displaystyle\int_1^{\expo{}} x(1 - \ln x)\,\mathrm{d}x$.
            \item \`A l'aide d'une seconde int\'egration par parties, calculer $J = \displaystyle\int_1^{\expo{}} f(x)\,\mathrm{d}x$.
            \item En d\'eduire l'aire, en cm$^2$, du domaine limit\'e par $(\mathscr{C})$, l'axe des abscisses et les droites
                  d'\'equations $x = 1$ et $x = \expo{}$.
            \end{enumerate}
\end{enumerate}
\begin{center}
On donne : $\expo{} \approx 2{,}72$ ; $\frac{4}{\expo{}} \approx 1{,}47$ ; $\expo{2} \approx 7{,}39$.
\end{center}

%%%%%%%%%%%%%%%%%%%%%%%%%%%%%%%%%%%%%%%%%%%%%%%%%%%%%%%%%%%%%%%%%%%%%%%%%%%%%%%%%%%%%%%%%%%%%%%%%%%%%%%%
\newpage
\section{Sujet D5}
\noindent \underline{\textsc{Exercice I}} (5 points) \newline
Le plan complexe est muni d'un rep\`ere orthonorm\'e direct $(O ; \vec{u}, \vec{v})$ (unit\'e 1~cm).
\begin{enumerate}
      \item Calculer $(3 - i)^2$. R\'esoudre dans $\mathbb{C}$ l'\'equation $z^2 - (1 + 3i)\,z - 4 + 3i = 0$.
      \item On note $A$ et $B$ les points d'affixes $z_A = 2 + i$ et $z_B = -1 + 2i$.
            Calculer $\frac{z_B}{z_A}$ et en d\'eduire la nature du triangle $OAB$.
      \item Soit $f$ l'application qui \`a tout point $M(z)$ associe le point $M'(z')$ tel que $z' = iz + 1 + 3i$.
            \begin{enumerate}
            \item Montrer que $f$ est une rotation dont on pr\'ecisera le centre et l'angle.
            \item Calculer l'affixe du point $C = f(A)$ et pr\'eciser la nature du triangle $BAC$.
            \end{enumerate}
      \item D\'eterminer et construire l'ensemble $(E)$ des points $M(z)$ tels que $|z - 2 - i| = |z + 1 - 2i|$.
            V\'erifier que $O$ appartient \`a $(E)$.
\end{enumerate}

\noindent \underline{\textsc{Exercice II}} (5 points) \newline
Une joueuse de basket r\'eussit chaque lancer franc avec la probabilit\'e $0{,}7$, ind\'ependamment des autres lancers.
Elle effectue cinq lancers francs. Soit $X$ le nombre de lancers r\'eussis.
\begin{enumerate}
      \item Quelle est la loi de $X$\,? Donner l'ensemble de ses valeurs.
      \item Calculer (\`a $10^{-4}$ pr\`es) la probabilit\'e qu'elle r\'eussisse exactement trois lancers, puis tous les lancers.
      \item Calculer la probabilit\'e qu'elle r\'eussisse au moins un lancer.
      \item Calculer l'esp\'erance et la variance de $X$. Chaque lancer r\'eussi rapporte 1 point :
            combien de points peut-elle esp\'erer marquer en moyenne\,?
      \item Combien de lancers doit-elle effectuer, au minimum, pour que la probabilit\'e d'en r\'eussir au moins un
            soit sup\'erieure ou \'egale \`a $0{,}999$\,?
\end{enumerate}

\noindent \underline{\textsc{Probl\`eme}} (10 points) \newline
Soit $f$ la fonction d\'efinie sur $\mathbb{R}$ par $f(x) = x - 2 + \dfrac{4}{\expo{x} + 1}$, et $(\mathscr{C})$ sa courbe
dans un rep\`ere orthonorm\'e (unit\'e 1~cm).
\begin{enumerate}
      \item Montrer que, pour tout r\'eel $x$, $\dfrac{4}{\expo{-x} + 1} = 4 - \dfrac{4}{\expo{x} + 1}$.
            En d\'eduire que $f$ est impaire. Interpr\'eter graphiquement.
      \item \begin{enumerate}
            \item Calculer les limites de $f$ en $+\infty$ et en $-\infty$.
            \item Montrer que les droites $(D_1) : y = x - 2$ et $(D_2) : y = x + 2$ sont asymptotes \`a $(\mathscr{C})$
                  respectivement en $+\infty$ et en $-\infty$. Pr\'eciser la position de $(\mathscr{C})$ par rapport \`a chacune.
            \end{enumerate}
      \item Montrer que $f'(x) = \dfrac{(\expo{x} - 1)^2}{(\expo{x} + 1)^2}$. Dresser le tableau de variations de $f$.
            Que peut-on dire de la tangente \`a $(\mathscr{C})$ en $O$\,?
      \item Montrer que l'\'equation $f(x) = 1$ admet une unique solution $\alpha$ et que $2 < \alpha < 3$.
      \item Tracer $(D_1)$, $(D_2)$ et $(\mathscr{C})$.
      \item \begin{enumerate}
            \item V\'erifier que, pour tout r\'eel $x$, $\dfrac{4}{\expo{x} + 1} = \dfrac{4\,\expo{-x}}{1 + \expo{-x}}$
                  et en d\'eduire une primitive de $x \mapsto \dfrac{4}{\expo{x} + 1}$.
            \item Soit $\lambda > 0$. Calculer l'aire $\mathscr{A}(\lambda)$ du domaine limit\'e par $(\mathscr{C})$, $(D_1)$ et
                  les droites d'\'equations $x = 0$ et $x = \lambda$, puis sa limite quand $\lambda$ tend vers $+\infty$.
            \end{enumerate}
      \item Justifier que $f$ admet une fonction r\'eciproque $f^{-1}$ d\'efinie sur $\mathbb{R}$. Est-elle d\'erivable en $0$\,?
\end{enumerate}
\begin{center}
On donne : $\expo{2} \approx 7{,}39$ ; $\expo{3} \approx 20{,}09$ ; $\ln 2 \approx 0{,}69$.
\end{center}

%%%%%%%%%%%%%%%%%%%%%%%%%%%%%%%%%%%%%%%%%%%%%%%%%%%%%%%%%%%%%%%%%%%%%%%%%%%%%%%%%%%%%%%%%%%%%%%%%%%%%%%%
\newpage
\section{Sujet D6}
\noindent \underline{\textsc{Exercice I}} (5 points) \newline
Le plan complexe est muni d'un rep\`ere orthonorm\'e direct $(O ; \vec{u}, \vec{v})$.
\begin{enumerate}
      \item Soit $\theta \in\ ]-\pi ; \pi[$ et $z = 1 + \expo{i\theta}$.
            \begin{enumerate}
            \item Montrer que $z = 2\cos\left(\frac{\theta}{2}\right) \expo{i\frac{\theta}{2}}$.
            \item En d\'eduire le module et un argument de $z$.
            \item Application : donner la forme exponentielle de $1 + \expo{i\frac{2\pi}{3}}$ et de $1 + i$.
            \end{enumerate}
      \item On consid\`ere les points $A(1)$, $B(i)$, $A'(1 + i)$ et $B'(-1 + i)$.
            \begin{enumerate}
            \item D\'eterminer l'\'ecriture complexe de la similitude directe $S$ telle que $S(A) = A'$ et $S(B) = B'$.
            \item Pr\'eciser son centre, son rapport et son angle.
            \end{enumerate}
      \item On pose $M_0 = A$ et, pour tout entier $n$, $M_{n+1} = S(M_n)$ ; on note $z_n$ l'affixe de $M_n$.
            \begin{enumerate}
            \item Montrer que $z_n = (1 + i)^n$ et l'\'ecrire sous forme exponentielle.
            \item Calculer $z_{12}$. Pour quelles valeurs de $n$ le point $M_n$ est-il sur la demi-droite $[O ; -\vec{u})$\,?
            \end{enumerate}
\end{enumerate}

\noindent \underline{\textsc{Exercice II}} (5 points) \newline
On a relev\'e la temp\'erature moyenne annuelle $y$ (en $^\circ$C) de six localit\'es en fonction de leur altitude $x$
(en centaines de m\`etres).
\begin{center}
\begin{tabular}{|c|c|c|c|c|c|c|}
\hline
Altitude $x_i$ (centaines de m) & 2 & 5 & 8 & 12 & 14 & 16 \\
\hline
Temp\'erature $y_i$ ($^\circ$C) & 27 & 25 & 23 & 20 & 19 & 17 \\
\hline
\end{tabular}
\end{center}
\begin{enumerate}
      \item Repr\'esenter le nuage de points et le point moyen $G$.
      \item Calculer $V(x)$, $V(y)$, la covariance $\sigma_{xy}$ et le coefficient de corr\'elation lin\'eaire $r$ (\`a $10^{-3}$ pr\`es).
            Que traduit le signe de $r$\,?
      \item D\'eterminer la droite de r\'egression de $y$ en $x$ (coefficients \`a $10^{-2}$ pr\`es).
            Interpr\'eter son coefficient directeur.
      \item Estimer la temp\'erature moyenne d'une localit\'e situ\'ee \`a 2\,000~m d'altitude.
      \item \`A l'aide de la droite de r\'egression de $x$ en $y$, estimer l'altitude d'une localit\'e
            de temp\'erature moyenne $15^\circ$C.
\end{enumerate}

\noindent \underline{\textsc{Probl\`eme}} (10 points) \newline
Soit $f$ la fonction d\'efinie sur $\mathbb{R}$ par $f(x) = \expo{2x} - 3\expo{x} + 2$, et $(\mathscr{C})$ sa courbe
dans un rep\`ere orthonorm\'e (unit\'e 2~cm).
\begin{enumerate}
      \item Calculer la limite de $f$ en $-\infty$ et interpr\'eter. Calculer les limites de $f(x)$ et de $\frac{f(x)}{x}$ en $+\infty$.
      \item Calculer $f'(x)$ et dresser le tableau de variations de $f$. On calculera la valeur exacte du minimum.
      \item V\'erifier que $f(x) = (\expo{x} - 1)(\expo{x} - 2)$ et en d\'eduire le signe de $f(x)$.
      \item \'Etudier la position de $(\mathscr{C})$ par rapport \`a son asymptote $(\Delta) : y = 2$.
      \item Donner une \'equation de la tangente $(T)$ \`a $(\mathscr{C})$ au point d'abscisse $0$. Tracer $(\Delta)$, $(T)$ et $(\mathscr{C})$.
      \item Calculer, en cm$^2$, l'aire du domaine limit\'e par $(\mathscr{C})$ et l'axe des abscisses entre les droites
            d'\'equations $x = 0$ et $x = \ln 2$.
      \item Soit $g$ la restriction de $f$ \`a l'intervalle $I = \left[\ln\frac{3}{2} ; +\infty\right[$.
            \begin{enumerate}
            \item Montrer que $g$ r\'ealise une bijection de $I$ sur un intervalle $J$ \`a pr\'eciser.
            \item D\'eterminer l'expression de $g^{-1}(y)$ pour $y \in J$.
            \end{enumerate}
\end{enumerate}
\begin{center}
On donne : $\ln 2 \approx 0{,}69$ ; $\ln 3 \approx 1{,}10$ ; $\ln\frac{3}{2} \approx 0{,}41$.
\end{center}

%%%%%%%%%%%%%%%%%%%%%%%%%%%%%%%%%%%%%%%%%%%%%%%%%%%%%%%%%%%%%%%%%%%%%%%%%%%%%%%%%%%%%%%%%%%%%%%%%%%%%%%%
%                                      CORRIGÉS (édition du professeur)
%%%%%%%%%%%%%%%%%%%%%%%%%%%%%%%%%%%%%%%%%%%%%%%%%%%%%%%%%%%%%%%%%%%%%%%%%%%%%%%%%%%%%%%%%%%%%%%%%%%%%%%%
\begin{corrige}
\newpage
\section{Corrig\'es des sujets d'entra\^inement de s\'erie D}
\subsection{Corrig\'e du sujet D1}
\noindent \underline{\textsc{Exercice I}}
\begin{enumerate}
\item \begin{enumerate}
      \item $P(i) = i^3 - (2+i)i^2 + (2+2i)i - 2i = -i + 2 + i + 2i - 2 - 2i = 0$.
      \item $(z - i)(z^2 + az + b) = z^3 + (a - i)z^2 + (b - ai)z - bi$. Par identification : $a - i = -2 - i$, $-bi = -2i$,
            d'o\`u $a = -2$, $b = 2$ ; on v\'erifie $b - ai = 2 + 2i$. Donc $P(z) = (z - i)(z^2 - 2z + 2)$.
      \item $z^2 - 2z + 2 = 0$ : $\Delta = -4 = (2i)^2$, $z = 1 \pm i$. Ainsi $\boxed{S = \{ i \,;\, 1 + i \,;\, 1 - i \}}$.
      \end{enumerate}
\item $Z = \dfrac{(1 - i) - (1 + i)}{i - (1 + i)} = \dfrac{-2i}{-1} = 2i = 2\expo{i\frac{\pi}{2}}$.
      Donc $\frac{BC}{BA} = |Z| = 2$ et $(\overrightarrow{BA}, \overrightarrow{BC}) = \frac{\pi}{2}$ : le triangle $ABC$ est
      rectangle en $B$, avec $BC = 2\,BA$.
\item \begin{enumerate}
      \item $z' - z_B = 2\expo{i\frac{\pi}{2}}(z - z_B)$, soit $z' = 2iz + (1+i)(1 - 2i) = \boxed{2iz + 3 - i}$.
            $S(A)$ : $2i \cdot i + 3 - i = 1 - i = z_C$ (c'est aussi la traduction de $Z = 2i$).
      \item Le centre $B$ de $(\Gamma)$ est fixe ; une similitude de rapport 2 multiplie les distances par 2 :
            l'image de $(\Gamma)$ est le cercle de centre $B$ et de rayon $4$.
      \end{enumerate}
\item $z_B = 1 + i = \sqrt{2}\left(\cos\frac{\pi}{4} + i \sin\frac{\pi}{4}\right)$. Alors
      $z_B^{2026} = (\sqrt{2})^{2026}\expo{i\frac{2026\pi}{4}} = 2^{1013}\expo{i\left(506\pi + \frac{\pi}{2}\right)} = \boxed{2^{1013}\, i}$.
\end{enumerate}

\noindent \underline{\textsc{Exercice II}} \newline
L'univers est l'ensemble des parties \`a 3 \'el\'ements d'un ensemble \`a 8 \'el\'ements : $\mathrm{C}_8^3 = 56$ tirages \'equiprobables.
\begin{enumerate}
\item $P(A) = \dfrac{\mathrm{C}_5^3 + \mathrm{C}_3^3}{56} = \dfrac{10 + 1}{56} = \dfrac{11}{56}$ ;
      $P(B) = 1 - P(\text{aucun vert}) = 1 - \dfrac{\mathrm{C}_5^3}{56} = 1 - \dfrac{10}{56} = \dfrac{23}{28}$.
\item $X$ prend les valeurs $0, 1, 2, 3$ et $P(X = k) = \dfrac{\mathrm{C}_3^k\,\mathrm{C}_5^{3-k}}{56}$ :
      \begin{center}
      \begin{tabular}{|c|c|c|c|c|}
      \hline
      $k$ & 0 & 1 & 2 & 3 \\ \hline
      $P(X = k)$ & $\frac{10}{56} = \frac{5}{28}$ & $\frac{30}{56} = \frac{15}{28}$ & $\frac{15}{56}$ & $\frac{1}{56}$ \\ \hline
      \end{tabular}
      \end{center}
      $E(X) = \dfrac{30 + 30 + 3}{56} = \dfrac{63}{56} = \dfrac{9}{8}$ ;
      $E(X^2) = \dfrac{30 + 60 + 9}{56} = \dfrac{99}{56}$, $V(X) = \dfrac{99}{56} - \dfrac{81}{64} = \dfrac{225}{448}$.
\item $G = 500X - 400$, donc $E(G) = 500 \times \frac{9}{8} - 400 = 162{,}5$~Ar $> 0$ : le jeu est favorable au joueur.
      Il est \'equitable si la mise vaut $500\,E(X) = 562{,}5$~Ar.
\item $P(X = 3) = \frac{1}{56}$ \`a chaque partie ; la probabilit\'e de ne jamais tirer trois verts en quatre parties
      ind\'ependantes est $\left(\frac{55}{56}\right)^4$. La probabilit\'e cherch\'ee est $1 - \left(\frac{55}{56}\right)^4 \approx 0{,}070$.
\end{enumerate}

\noindent \underline{\textsc{Probl\`eme}} \newline
\underline{Partie A}
\begin{enumerate}
\item $g'(x) = -\expo{-x} + (x+1)\expo{-x} = x\,\expo{-x}$, du signe de $x$ : $g$ est d\'ecroissante sur $]-\infty ; 0]$
      et croissante sur $[0 ; +\infty[$.
\item Le minimum de $g$ est $g(0) = 1 - 1 = 0$. Donc $g(x) \geq 0$, avec \'egalit\'e seulement en $0$ (stricte monotonie de part et d'autre).
\end{enumerate}
\underline{Partie B}
\begin{enumerate}
\item En $+\infty$ : $(x + 2)\expo{-x} = x\expo{-x} + 2\expo{-x} \to 0$ (croissances compar\'ees), donc $f(x) \to +\infty$.
      En $-\infty$ : $x + 2 \to -\infty$ et $\expo{-x} \to +\infty$, donc $(x+2)\expo{-x} \to -\infty$ et $f(x) \to -\infty$.
\item $f'(x) = 1 + \expo{-x} - (x + 2)\expo{-x} = 1 - (x+1)\expo{-x} = g(x)$. Comme $g \geq 0$ et ne s'annule qu'en $0$,
      $f$ est strictement croissante sur $\mathbb{R}$ : elle cro\^it de $-\infty$ \`a $+\infty$ en passant par $f(0) = 1$.
\item \begin{enumerate}
      \item $f(x) - (x - 1) = (x + 2)\expo{-x} \to 0$ en $+\infty$ : $(D)$ est asymptote. Cette diff\'erence a le signe de $x + 2$ :
            $(\mathscr{C})$ est au-dessus de $(D)$ sur $]-2 ; +\infty[$, au-dessous sur $]-\infty ; -2[$ ; ils se coupent en $(-2 ; -3)$.
      \item $\frac{f(x)}{x} = 1 - \frac{1}{x} + \left(1 + \frac{2}{x}\right)\expo{-x} \to +\infty$ en $-\infty$ :
            $(\mathscr{C})$ admet une branche parabolique de direction $(Oy)$ en $-\infty$.
      \end{enumerate}
\item $f''(x) = g'(x) = x\,\expo{-x}$ s'annule en changeant de signe en $0$ : $I(0 ; f(0)) = (0 ; 1)$ est un point d'inflexion.
      Comme $f'(0) = g(0) = 0$, la tangente en $I$ est horizontale : $y = 1$.
\item $f$ est continue et strictement croissante sur $\mathbb{R}$, de $-\infty$ \`a $+\infty$ : elle s'annule une seule fois.
      $f(-1{,}5) = -2{,}5 + 0{,}5\,\expo{1{,}5} \approx -0{,}26 < 0$ et $f(-1) = -2 + \expo{} \approx 0{,}72 > 0$, donc
      $-1{,}5 < \alpha < -1$ ($\alpha \approx -1{,}41$).
\item Trac\'e : asymptote $(D)$, point $(-2 ; -3)$, tangente horizontale en $I(0;1)$, $(\mathscr{C})$ coupe $(Ox)$ en $\alpha$.
\item Sur $[0 ; \lambda]$, $(\mathscr{C})$ est au-dessus de $(D)$, donc
      $\mathscr{A}(\lambda) = \displaystyle\int_0^\lambda (x+2)\expo{-x}\,\mathrm{d}x$ (en unit\'es d'aire, ici $1$~cm$^2$).
      Avec $u = x + 2$, $v' = \expo{-x}$ :
      $\mathscr{A}(\lambda) = \Big[-(x+2)\expo{-x}\Big]_0^\lambda + \displaystyle\int_0^\lambda \expo{-x}\,\mathrm{d}x
      = -(\lambda + 2)\expo{-\lambda} + 2 + 1 - \expo{-\lambda} = \boxed{3 - (\lambda + 3)\expo{-\lambda}}$~cm$^2$,
      et $\lim\limits_{\lambda \to +\infty} \mathscr{A}(\lambda) = 3$~cm$^2$.
\item $f$ est continue et strictement croissante de $\mathbb{R}$ sur $\mathbb{R}$ : c'est une bijection, elle admet une r\'eciproque
      $f^{-1}$ d\'efinie sur $\mathbb{R}$. Comme $f(0) = 1$ et $f'(0) = 0$, $f^{-1}$ n'est pas d\'erivable en $1$ : la courbe de $f^{-1}$
      (sym\'etrique de $(\mathscr{C})$ par rapport \`a la droite $y = x$) a une tangente verticale au point $(1 ; 0)$.
\end{enumerate}

%%%%%%%%%%%%%%%%%%%%%%%%%%%%%%%%%%%%%%%%%%%%%%%%%%%%%%%%%%%%%%%%%%%%%%%%%%%%%%%%%%%%%%%%%%%%%%%%%%%%%%%%
\subsection{Corrig\'e du sujet D2}
\noindent \underline{\textsc{Exercice I}}
\begin{enumerate}
\item $\Delta = 12 - 16 = -4 = (2i)^2$ : $z_A = \sqrt{3} + i$ et $z_B = \sqrt{3} - i$.
\item $|z_A| = 2$ et $z_A = 2\left(\frac{\sqrt{3}}{2} + \frac{1}{2}i\right) = 2\expo{i\frac{\pi}{6}}$ ; $z_B = \overline{z_A} = 2\expo{-i\frac{\pi}{6}}$.
      $z_A^6 = 2^6\expo{i\pi} = -64$. $\frac{z_A}{z_B} = \expo{i\frac{\pi}{6} + i\frac{\pi}{6}} = \expo{i\frac{\pi}{3}}$ :
      $OA = OB$ et $(\overrightarrow{OB}, \overrightarrow{OA}) = \frac{\pi}{3}$, donc $OAB$ est \'equilat\'eral.
\item $z_{A'} = \expo{i\frac{\pi}{3}} z_A = 2\expo{i\frac{\pi}{2}} = 2i$.
\item \begin{enumerate}
      \item $1 + i\sqrt{3} = 2\expo{i\frac{\pi}{3}}$ : $S$ est la similitude directe de centre $O$ (seul point fixe),
            de rapport $2$ et d'angle $\frac{\pi}{3}$.
      \item $S(A)$ : $(1 + i\sqrt{3})(\sqrt{3} + i) = 4i$ ; $S(B)$ : $(1 + i\sqrt{3})(\sqrt{3} - i) = 2\sqrt{3} + 2i = 2 z_A$.
            Comme $z_{S(B)} = 2 z_A$, les vecteurs $\overrightarrow{OS(B)}$ et $\overrightarrow{OA}$ sont colin\'eaires :
            $O$, $A$, $S(B)$ sont align\'es.
      \end{enumerate}
\item $|z - 2i| = |z - \sqrt{3} - i|$ signifie $MA' = MA$ : $(E)$ est la m\'ediatrice de $[AA']$.
      Comme $OA = OA' = 2$, $O \in (E)$ ; le milieu de $[AA']$ a pour affixe $\frac{\sqrt{3}}{2} + \frac{3}{2}i$,
      donc $(E)$ est la droite d'\'equation $y = \sqrt{3}\,x$.
\end{enumerate}

\noindent \underline{\textsc{Exercice II}} \newline
$n = 6$, $\sum x_i = 21$, $\sum y_i = 450$, $\sum x_i^2 = 91$, $\sum y_i^2 = 35\,290$, $\sum x_i y_i = 1\,739$.
\begin{enumerate}
\item Les points sont presque align\'es sur une droite montante.
\item $G(3{,}5 ; 75)$.
\item $V(x) = \frac{91}{6} - 3{,}5^2 = \frac{35}{12} \approx 2{,}917$ ; $\sigma_{xy} = \frac{1\,739}{6} - 3{,}5 \times 75 = \frac{82}{3} \approx 27{,}333$ ;
      $V(y) = \frac{35\,290}{6} - 75^2 = \frac{770}{3} \approx 256{,}667$ ; $r = \frac{\sigma_{xy}}{\sqrt{V(x)V(y)}} \approx 0{,}999$.
      $r$ est tr\`es proche de $1$ : l'ajustement affine est justifi\'e.
\item $a = \frac{\sigma_{xy}}{V(x)} = \frac{328}{35} \approx 9{,}37$ et $b = 75 - \frac{328}{35} \times 3{,}5 = 42{,}2$ :
      $\boxed{y = 9{,}37\,x + 42{,}2}$.
\item \begin{enumerate}
      \item 2027 correspond \`a $x = 8$ : $y \approx 9{,}37 \times 8 + 42{,}2 \approx 117{,}2$, soit environ $11\,700$ abonn\'es.
      \item $9{,}37\,x + 42{,}2 > 150 \iff x > \frac{107{,}8}{9{,}37} \approx 11{,}5$, donc $x \geq 12$ : \`a partir de 2031.
      \end{enumerate}
\end{enumerate}

\noindent \underline{\textsc{Probl\`eme}} \newline
\underline{Partie A}
\begin{enumerate}
\item $g'(x) = 2x + \frac{1}{x} > 0$ sur $]0 ; +\infty[$ : $g$ est strictement croissante. $\lim_{0^+} g = -\infty$
      (car $\ln x \to -\infty$) et $\lim_{+\infty} g = +\infty$.
\item $g$ est continue strictement croissante de $]0;+\infty[$ sur $\mathbb{R}$ : $g(x) = 0$ a une unique solution $\alpha$.
      $g(0{,}6) \approx 0{,}36 - 0{,}51 = -0{,}15 < 0$ et $g(0{,}7) \approx 0{,}49 - 0{,}36 = 0{,}13 > 0$, d'o\`u $0{,}6 < \alpha < 0{,}7$.
\item $g(x) < 0$ sur $]0 ; \alpha[$, $g(\alpha) = 0$, $g(x) > 0$ sur $]\alpha ; +\infty[$.
\end{enumerate}
\underline{Partie B}
\begin{enumerate}
\item En $0^+$ : $1 + \ln x \to -\infty$ et $\frac{1}{x} \to +\infty$, donc $\frac{1 + \ln x}{x} \to -\infty$ et $f(x) \to +\infty$ :
      l'axe des ordonn\'ees est asymptote verticale. En $+\infty$ : $\frac{1}{x} \to 0$ et $\frac{\ln x}{x} \to 0$, donc $f(x) \to +\infty$.
\item $\left(\frac{1 + \ln x}{x}\right)' = \frac{\frac{1}{x}\cdot x - (1 + \ln x)}{x^2} = \frac{-\ln x}{x^2}$, donc
      $f'(x) = 1 + \frac{\ln x}{x^2} = \frac{x^2 + \ln x}{x^2} = \frac{g(x)}{x^2}$.
      $f$ est d\'ecroissante sur $]0 ; \alpha]$ et croissante sur $[\alpha ; +\infty[$, de minimum $f(\alpha)$.
\item $g(\alpha) = 0$ donne $\ln\alpha = -\alpha^2$, donc $f(\alpha) = \alpha - \frac{1 - \alpha^2}{\alpha} = \frac{2\alpha^2 - 1}{\alpha}$
      ($\approx -0{,}23$).
\item \begin{enumerate}
      \item $f(x) - x = -\frac{1 + \ln x}{x} \to 0$ en $+\infty$ : $(D)$ est asymptote oblique.
      \item $f(x) - x$ a le signe de $-(1 + \ln x)$ : $(\mathscr{C})$ est au-dessus de $(D)$ sur $\left]0 ; \frac{1}{\expo{}}\right[$,
            au-dessous sur $\left]\frac{1}{\expo{}} ; +\infty\right[$ ; elles se coupent au point $\left(\frac{1}{\expo{}} ; \frac{1}{\expo{}}\right)$.
      \end{enumerate}
\item $f(1) = 1 - \frac{1 + 0}{1} = 0$. Sur $[\alpha ; +\infty[$, $f$ est strictement croissante de $f(\alpha) < 0$ \`a $+\infty$ :
      une seule solution, qui est $1$. Sur $]0 ; \alpha]$, $f$ est strictement d\'ecroissante de $+\infty$ \`a $f(\alpha) < 0$ : une seule
      solution $\beta$. $f(0{,}4) \approx 0{,}4 - \frac{1 - 0{,}92}{0{,}4} = 0{,}2 > 0$ et $f(0{,}5) \approx 0{,}5 - \frac{1 - 0{,}69}{0{,}5} = -0{,}12 < 0$,
      donc $0{,}4 < \beta < 0{,}5$.
\item Trac\'e : asymptotes $x = 0$ et $(D)$, minimum en $\alpha$, points d'abscisses $\beta$ et $1$ sur l'axe des abscisses.
\item Sur $\left[\frac{1}{\expo{}} ; \expo{}\right]$, $(\mathscr{C})$ est au-dessous de $(D)$ ; l'aire vaut, en unit\'es d'aire,
      $\displaystyle\int_{1/\expo{}}^{\expo{}} \frac{1 + \ln x}{x}\,\mathrm{d}x = \left[\frac{(1 + \ln x)^2}{2}\right]_{1/\expo{}}^{\expo{}} = \frac{4}{2} - 0 = 2$.
      L'unit\'e d'aire valant $4$~cm$^2$, l'aire est $\boxed{8 \text{ cm}^2}$.
\end{enumerate}

%%%%%%%%%%%%%%%%%%%%%%%%%%%%%%%%%%%%%%%%%%%%%%%%%%%%%%%%%%%%%%%%%%%%%%%%%%%%%%%%%%%%%%%%%%%%%%%%%%%%%%%%
\subsection{Corrig\'e du sujet D3}
\noindent \underline{\textsc{Exercice I}} \newline
On pose $q = \frac{1 + i}{2}$.
\begin{enumerate}
\item $z_1 = 2 + 2i$, $z_2 = (1 + i)^2 = 2i$, $z_3 = i(1 + i) = -1 + i$, $z_4 = \frac{(1+i)(-1+i)}{2} = -1$.
\item $q = \frac{\sqrt{2}}{2}\expo{i\frac{\pi}{4}}$. Par r\'ecurrence imm\'ediate, $z_n = q^n z_0 = 4\left(\frac{\sqrt{2}}{2}\right)^n \expo{i\frac{n\pi}{4}}$.
\item $r_{n+1} = |q|\,r_n = \frac{\sqrt{2}}{2}\, r_n$ : suite g\'eom\'etrique de raison $\frac{\sqrt{2}}{2} \in\ ]0 ; 1[$ et de premier terme 4,
      donc $\lim r_n = 0$.
\item $\dfrac{z_n - z_{n+1}}{-z_{n+1}} = \dfrac{(1 - q)z_n}{-q z_n} = 1 - \dfrac{1}{q} = 1 - (1 - i) = i$
      (car $\frac{1}{q} = \frac{2}{1 + i} = 1 - i$). Donc $M_{n+1}M_n = M_{n+1}O$ et $(\overrightarrow{M_{n+1}O}, \overrightarrow{M_{n+1}M_n}) = \frac{\pi}{2}$ :
      le triangle $OM_{n+1}M_n$ est rectangle isoc\`ele en $M_{n+1}$.
\item $\ell_n = |z_{n+1} - z_n| = |q - 1|\, r_n = \frac{\sqrt{2}}{2} r_n = 2\sqrt{2}\left(\frac{\sqrt{2}}{2}\right)^n$.
      Somme de termes d'une suite g\'eom\'etrique : $L_n = 2\sqrt{2}\, \dfrac{1 - \left(\frac{\sqrt{2}}{2}\right)^n}{1 - \frac{\sqrt{2}}{2}}$,
      et $\lim L_n = \dfrac{2\sqrt{2}}{1 - \frac{\sqrt{2}}{2}} = \dfrac{4\sqrt{2}}{2 - \sqrt{2}} = \boxed{4 + 4\sqrt{2}} \approx 9{,}66$ (unit\'es de longueur).
\item $M_n \in [O ; \vec{u})$ si et seulement si $\frac{n\pi}{4}$ est un multiple de $2\pi$, c'est-\`a-dire si $n$ est un multiple de $8$
      (par exemple $z_8 = \frac{1}{4}$).
\end{enumerate}

\noindent \underline{\textsc{Exercice II}} \newline
On note $D$ l'\'ev\`enement << la mangue est ab\^im\'ee >>.
\begin{enumerate}
\item Arbre : trois branches $P_1$ ($0{,}5$), $P_2$ ($0{,}3$), $P_3$ ($0{,}2$), puis de chacune $D$ et $\overline{D}$ avec
      $P_{P_1}(D) = 0{,}04$, $P_{P_2}(D) = 0{,}05$, $P_{P_3}(D) = 0{,}1$.
\item Formule des probabilit\'es totales : $P(D) = 0{,}5 \times 0{,}04 + 0{,}3 \times 0{,}05 + 0{,}2 \times 0{,}1 = 0{,}02 + 0{,}015 + 0{,}02 = 0{,}055$.
\item $P_D(P_3) = \dfrac{P(P_3 \cap D)}{P(D)} = \dfrac{0{,}02}{0{,}055} = \dfrac{4}{11} \approx 0{,}364$.
\item \begin{enumerate}
      \item Trois \'epreuves ind\'ependantes identiques \`a deux issues : $X$ suit la loi binomiale $\mathscr{B}(3 ; 0{,}055)$.
            $P(X = 0) = 0{,}945^3 \approx 0{,}8439$ et $P(X \geq 1) = 1 - 0{,}945^3 \approx 0{,}1561$.
      \item $E(X) = 3 \times 0{,}055 = 0{,}165$.
      \end{enumerate}
\end{enumerate}

\noindent \underline{\textsc{Probl\`eme}}
\begin{enumerate}
\item En $+\infty$ : $f(x) = x^2\expo{-x} + 2x\expo{-x} + \expo{-x} \to 0$ (croissances compar\'ees) : l'axe des abscisses est asymptote en $+\infty$.
      En $-\infty$ : $(x+1)^2 \to +\infty$ et $\expo{-x} \to +\infty$, donc $f(x) \to +\infty$ ;
      $\frac{f(x)}{x} = \frac{(x+1)^2}{x}\expo{-x} \to -\infty$ : branche parabolique de direction $(Oy)$.
\item $f'(x) = 2(x+1)\expo{-x} - (x+1)^2\expo{-x} = (x+1)(2 - x - 1)\expo{-x} = (x+1)(1-x)\expo{-x}$.
      $f$ est d\'ecroissante sur $]-\infty ; -1]$, croissante sur $[-1 ; 1]$, d\'ecroissante sur $[1 ; +\infty[$ ;
      $f(-1) = 0$ (minimum) et $f(1) = \frac{4}{\expo{}} \approx 1{,}47$ (maximum).
\item \begin{enumerate}
      \item $f(0) = 1$, $f'(0) = 1$ : $(T) : y = x + 1$.
      \item $f(x) - (x+1) = (x+1)\big[(x+1)\expo{-x} - 1\big] = (x+1)(x + 1 - \expo{x})\expo{-x}$.
      \item $x + 1 - \expo{x} \leq 0$ (nul seulement en $0$) et $\expo{-x} > 0$, donc $f(x) - (x + 1)$ a le signe de $-(x+1)$ :
            $(\mathscr{C})$ est au-dessus de $(T)$ sur $]-\infty ; -1[$ et au-dessous sur $]-1 ; +\infty[$
            (avec contact au point $(0 ; 1)$) ; ils se coupent aussi en $(-1 ; 0)$.
      \end{enumerate}
\item Sur $]-\infty ; -1]$, $f$ d\'ecro\^it strictement de $+\infty$ \`a $0$ : une solution $\gamma$ ; $f(-2) = \expo{2} > 1$ et $f(-1) = 0 < 1$,
      donc $\gamma \in\ ]-2 ; -1[$. Sur $[-1 ; 1]$, $f$ cro\^it strictement de $0$ \`a $\frac{4}{\expo{}}$ : une solution, et $f(0) = 1$, c'est $0$.
      Sur $[1 ; +\infty[$, $f$ d\'ecro\^it strictement de $\frac{4}{\expo{}} > 1$ vers $0$ : une solution $\alpha$ ;
      $f(2) = \frac{9}{\expo{2}} \approx 1{,}22 > 1$ et $f(3) = \frac{16}{\expo{3}} \approx 0{,}80 < 1$, donc $\alpha \in\ ]2 ; 3[$.
\item Trac\'e : minimum $(-1 ; 0)$ (tangente horizontale), maximum $\left(1 ; \frac{4}{\expo{}}\right)$, tangente $(T)$, asymptote $y = 0$ en $+\infty$.
\item \begin{enumerate}
      \item $F'(x) = -(2x + 4)\expo{-x} + (x^2 + 4x + 5)\expo{-x} = (x^2 + 2x + 1)\expo{-x} = f(x)$.
      \item $f \geq 0$, donc l'aire vaut $F(1) - F(0) = -\frac{10}{\expo{}} + 5 = 5 - \frac{10}{\expo{}}$ unit\'es d'aire, soit
            $\boxed{4\left(5 - \frac{10}{\expo{}}\right) = 20 - \frac{40}{\expo{}} \approx 5{,}28 \text{ cm}^2}$.
      \end{enumerate}
\item $f$ est d\'ecroissante sur $[n ; n+1] \subset [1 ; +\infty[$ : pour $x \in [n ; n+1]$, $f(n+1) \leq f(x) \leq f(n)$ ; en int\'egrant
      sur un intervalle de longueur 1, $f(n+1) \leq u_n \leq f(n)$. Comme $\lim_{+\infty} f = 0$, le th\'eor\`eme des gendarmes donne $\lim u_n = 0$.
\end{enumerate}

%%%%%%%%%%%%%%%%%%%%%%%%%%%%%%%%%%%%%%%%%%%%%%%%%%%%%%%%%%%%%%%%%%%%%%%%%%%%%%%%%%%%%%%%%%%%%%%%%%%%%%%%
\subsection{Corrig\'e du sujet D4}
\noindent \underline{\textsc{Exercice I}}
\begin{enumerate}
\item $(2 + i)^2 = 4 + 4i - 1 = 3 + 4i$. $\Delta = (4 + 3i)^2 - 4(1 + 5i) = 7 + 24i - 4 - 20i = 3 + 4i = (2+i)^2$.
      Solutions : $\frac{4 + 3i + 2 + i}{2} = 3 + 2i$ et $\frac{4 + 3i - 2 - i}{2} = 1 + i$.
\item \begin{enumerate}
      \item $z_C - z_B = -i(z_A - z_B) = -i(-2 - i) = -1 + 2i$, donc $z_C = 2 + 4i$.
            $z_D - z_A = i(z_B - z_A) = i(2 + i) = -1 + 2i$, donc $z_D = 3i$.
      \item $\overrightarrow{AD}$ et $\overrightarrow{BC}$ ont la m\^eme affixe $-1 + 2i$ : $ABCD$ est un parall\'elogramme ;
            $AD = AB = \sqrt{5}$ et $(\overrightarrow{AB}, \overrightarrow{AD}) = \frac{\pi}{2}$ : c'est un carr\'e (direct).
            Son centre est le milieu de $[AC]$ : $z_\Omega = \frac{3}{2} + \frac{5}{2}i$.
      \end{enumerate}
\item \begin{enumerate}
      \item $s : z' - z_A = a(z - z_A)$ avec $a = \frac{z_C - z_A}{z_B - z_A} = \frac{1 + 3i}{2 + i} = \frac{(1 + 3i)(2 - i)}{5} = 1 + i$.
            Donc $z' = (1 + i)z + 1 - i$ ; rapport $|1 + i| = \sqrt{2}$, angle $\frac{\pi}{4}$.
      \item $s(C)$ : $(1+i)(2 + 4i) + 1 - i = -2 + 6i + 1 - i = -1 + 5i$ ; $s(D)$ : $(1 + i)(3i) + 1 - i = -3 + 3i + 1 - i = -2 + 2i$.
      \item Une similitude de rapport $\sqrt{2}$ multiplie les aires par $2$ : l'aire de l'image est $2 \times AB^2 = 2 \times 5 = 10$~cm$^2$.
      \end{enumerate}
\end{enumerate}

\noindent \underline{\textsc{Exercice II}}
\begin{description}
\item[I.] \begin{enumerate}
      \item Quatre \'epreuves ind\'ependantes \`a deux issues, de m\^eme probabilit\'e de succ\`es $\frac{1}{3}$ : $X \sim \mathscr{B}\left(4 ; \frac{1}{3}\right)$.
      \item $P(X \geq 3) = \mathrm{C}_4^3\left(\frac{1}{3}\right)^3\frac{2}{3} + \left(\frac{1}{3}\right)^4 = \frac{8}{81} + \frac{1}{81} = \frac{1}{9}$ ;
            $P(X \geq 1) = 1 - \left(\frac{2}{3}\right)^4 = \frac{65}{81}$.
      \item $N = 3X - (4 - X) = 4X - 4$, donc $E(N) = 4 \times \frac{4}{3} - 4 = \frac{4}{3}$.
      \end{enumerate}
\item[II.] $n = 6$, $\sum x_i = 36$, $\sum y_i = 65$, $\sum x_i^2 = 258$, $\sum y_i^2 = 743$, $\sum x_i y_i = 430$.
      \begin{enumerate}
      \item $G\left(6 ; \frac{65}{6}\right) \approx (6 ; 10{,}83)$. $V(x) = 43 - 36 = 7$, $\sigma_{xy} = \frac{430}{6} - 65 = \frac{20}{3} \approx 6{,}67$,
            $V(y) = \frac{743}{6} - \left(\frac{65}{6}\right)^2 \approx 6{,}47$ et $r \approx \frac{6{,}67}{\sqrt{7 \times 6{,}47}} \approx 0{,}99$.
      \item $a = \frac{20/3}{7} = \frac{20}{21} \approx 0{,}95$, $b = \frac{65}{6} - \frac{20}{21} \times 6 = \frac{215}{42} \approx 5{,}12$ :
            $y = 0{,}95\,x + 5{,}12$.
      \item Pour $x = 9$ : $y \approx 0{,}95 \times 9 + 5{,}12 \approx 13{,}7$.
      \end{enumerate}
\end{description}

\noindent \underline{\textsc{Probl\`eme}}
\begin{enumerate}
\item \begin{enumerate}
      \item $f(x) = x - 2x\ln x + x(\ln x)^2 \to 0 = f(0)$ en $0^+$ : $f$ est continue en $0$.
      \item $\frac{f(x) - f(0)}{x} = (1 - \ln x)^2 \to +\infty$ en $0^+$ : $f$ n'est pas d\'erivable en $0$ ;
            $(\mathscr{C})$ admet en $O$ une demi-tangente verticale.
      \end{enumerate}
\item $\lim_{+\infty} f = +\infty$ et $\frac{f(x)}{x} = (1 - \ln x)^2 \to +\infty$ : branche parabolique de direction $(Oy)$.
\item $f'(x) = (1 - \ln x)^2 + x \cdot 2(1 - \ln x)\left(-\frac{1}{x}\right) = (1 - \ln x)(1 - \ln x - 2) = (\ln x - 1)(\ln x + 1)$.
      $f'(x) > 0$ si $\ln x < -1$ ou $\ln x > 1$ : $f$ est croissante sur $\left[0 ; \frac{1}{\expo{}}\right]$, d\'ecroissante sur
      $\left[\frac{1}{\expo{}} ; \expo{}\right]$, croissante sur $[\expo{} ; +\infty[$ ;
      $f\left(\frac{1}{\expo{}}\right) = \frac{1}{\expo{}}(1 + 1)^2 = \frac{4}{\expo{}}$ et $f(\expo{}) = 0$.
\item $f(1) = 1$, $f'(1) = (0 - 1)(0 + 1) = -1$ : $(T) : y = -x + 2$.
\item $f(1) = 1 \times 1^2 = 1$. Sur $\left]0 ; \frac{1}{\expo{}}\right]$, $f$ cro\^it strictement de $0$ \`a $\frac{4}{\expo{}} > 1$ : une solution.
      Sur $\left[\frac{1}{\expo{}} ; \expo{}\right]$, $f$ d\'ecro\^it strictement de $\frac{4}{\expo{}}$ \`a $0$ : une solution, qui est $1$.
      Sur $[\expo{} ; +\infty[$, $f$ cro\^it strictement de $0$ \`a $+\infty$ : une solution. Au total, trois solutions
      (environ $0{,}08$ ; $1$ ; $4{,}38$).
\item Trac\'e : demi-tangente verticale en $O$, maximum local $\left(\frac{1}{\expo{}} ; \frac{4}{\expo{}}\right)$,
      $(\mathscr{C})$ tangente \`a l'axe des abscisses en $(\expo{} ; 0)$, tangente $(T)$ en $(1 ; 1)$.
\item \begin{enumerate}
      \item $u = 1 - \ln x$, $v' = x$ : $K = \left[\frac{x^2}{2}(1 - \ln x)\right]_1^{\expo{}} + \displaystyle\int_1^{\expo{}} \frac{x^2}{2}\cdot\frac{1}{x}\,\mathrm{d}x
            = -\frac{1}{2} + \left[\frac{x^2}{4}\right]_1^{\expo{}} = \frac{\expo{2} - 3}{4}$.
      \item $u = (1 - \ln x)^2$, $v' = x$, $u' = -\frac{2(1 - \ln x)}{x}$ :
            $J = \left[\frac{x^2}{2}(1 - \ln x)^2\right]_1^{\expo{}} + \displaystyle\int_1^{\expo{}} x(1 - \ln x)\,\mathrm{d}x = -\frac{1}{2} + K = \frac{\expo{2} - 5}{4}$.
      \item $f \geq 0$, donc l'aire est $J$ unit\'es d'aire, soit $4J = \boxed{(\expo{2} - 5) \text{ cm}^2} \approx 2{,}39$~cm$^2$.
      \end{enumerate}
\end{enumerate}

%%%%%%%%%%%%%%%%%%%%%%%%%%%%%%%%%%%%%%%%%%%%%%%%%%%%%%%%%%%%%%%%%%%%%%%%%%%%%%%%%%%%%%%%%%%%%%%%%%%%%%%%
\subsection{Corrig\'e du sujet D5}
\noindent \underline{\textsc{Exercice I}}
\begin{enumerate}
\item $(3 - i)^2 = 8 - 6i$. $\Delta = (1 + 3i)^2 + 4(4 - 3i) = -8 + 6i + 16 - 12i = 8 - 6i = (3 - i)^2$.
      Solutions : $\frac{1 + 3i + 3 - i}{2} = 2 + i$ et $\frac{1 + 3i - 3 + i}{2} = -1 + 2i$.
\item $\frac{z_B}{z_A} = \frac{(-1 + 2i)(2 - i)}{5} = \frac{5i}{5} = i$ : $OB = OA$ et $(\overrightarrow{OA}, \overrightarrow{OB}) = \frac{\pi}{2}$ :
      $OAB$ est rectangle isoc\`ele en $O$.
\item \begin{enumerate}
      \item $z' = iz + 1 + 3i$ avec $|i| = 1$, $\arg i = \frac{\pi}{2}$ : $f$ est une rotation d'angle $\frac{\pi}{2}$ ; son centre $\omega$
            v\'erifie $\omega = i\omega + 1 + 3i$, soit $\omega = \frac{1 + 3i}{1 - i} = \frac{(1+3i)(1+i)}{2} = -1 + 2i$ : le centre est $B$.
      \item $z_C = i(2 + i) + 1 + 3i = 5i$. $C$ est l'image de $A$ par la rotation de centre $B$ et d'angle $\frac{\pi}{2}$ :
            $BAC$ est rectangle isoc\`ele en $B$.
      \end{enumerate}
\item $|z - z_A| = |z - z_B|$ : $(E)$ est la m\'ediatrice de $[AB]$. $|z_A| = |z_B| = \sqrt{5}$, donc $O \in (E)$ ;
      le milieu de $[AB]$ a pour affixe $\frac{1}{2} + \frac{3}{2}i$ : $(E)$ est la droite d'\'equation $y = 3x$.
\end{enumerate}

\noindent \underline{\textsc{Exercice II}}
\begin{enumerate}
\item Cinq \'epreuves ind\'ependantes \`a deux issues, $p = 0{,}7$ : $X \sim \mathscr{B}(5 ; 0{,}7)$, valeurs $0, 1, \ldots, 5$.
\item $P(X = 3) = \mathrm{C}_5^3\, 0{,}7^3\, 0{,}3^2 = 10 \times 0{,}343 \times 0{,}09 = 0{,}3087$ ; $P(X = 5) = 0{,}7^5 = 0{,}16807 \approx 0{,}1681$.
\item $P(X \geq 1) = 1 - 0{,}3^5 = 0{,}99757$.
\item $E(X) = 5 \times 0{,}7 = 3{,}5$ et $V(X) = 5 \times 0{,}7 \times 0{,}3 = 1{,}05$ : en moyenne $3{,}5$ points.
\item Avec $n$ lancers, $P(\text{au moins un}) = 1 - 0{,}3^n \geq 0{,}999 \iff 0{,}3^n \leq 0{,}001 \iff n \geq \frac{\ln 0{,}001}{\ln 0{,}3} \approx 5{,}74$
      (on divise par $\ln 0{,}3 < 0$) : il faut au moins $6$ lancers.
\end{enumerate}

\noindent \underline{\textsc{Probl\`eme}}
\begin{enumerate}
\item $\frac{4}{\expo{-x} + 1} = \frac{4\expo{x}}{1 + \expo{x}} = \frac{4(\expo{x} + 1) - 4}{\expo{x} + 1} = 4 - \frac{4}{\expo{x} + 1}$.
      Alors $f(-x) = -x - 2 + 4 - \frac{4}{\expo{x} + 1} = -\left(x - 2 + \frac{4}{\expo{x} + 1}\right) = -f(x)$ : $f$ est impaire,
      $(\mathscr{C})$ est sym\'etrique par rapport \`a $O$.
\item \begin{enumerate}
      \item En $+\infty$, $\frac{4}{\expo{x} + 1} \to 0$ et $f(x) \to +\infty$ ; en $-\infty$, $\frac{4}{\expo{x} + 1} \to 4$ et $f(x) \to -\infty$.
      \item $f(x) - (x - 2) = \frac{4}{\expo{x} + 1} \to 0$ en $+\infty$ et cette diff\'erence est $> 0$ : $(D_1)$ est asymptote,
            $(\mathscr{C})$ est au-dessus de $(D_1)$. $f(x) - (x + 2) = \frac{4}{\expo{x} + 1} - 4 = -\frac{4\expo{x}}{\expo{x} + 1} \to 0$
            en $-\infty$ et est $< 0$ : $(D_2)$ est asymptote, $(\mathscr{C})$ est au-dessous de $(D_2)$.
      \end{enumerate}
\item $f'(x) = 1 - \frac{4\expo{x}}{(\expo{x} + 1)^2} = \frac{\expo{2x} + 2\expo{x} + 1 - 4\expo{x}}{(\expo{x} + 1)^2} = \frac{(\expo{x} - 1)^2}{(\expo{x} + 1)^2} \geq 0$,
      nul seulement en $0$ : $f$ est strictement croissante de $-\infty$ \`a $+\infty$, $f(0) = 0$. La tangente en $O$ est horizontale ;
      comme $O$ est centre de sym\'etrie, $(\mathscr{C})$ traverse sa tangente en $O$ (point d'inflexion).
\item $f$ est continue et strictement croissante de $\mathbb{R}$ sur $\mathbb{R}$ : $f(x) = 1$ a une unique solution $\alpha$.
      $f(2) = \frac{4}{\expo{2} + 1} \approx 0{,}48 < 1$ et $f(3) = 1 + \frac{4}{\expo{3} + 1} \approx 1{,}19 > 1$, donc $2 < \alpha < 3$.
\item Trac\'e : les deux asymptotes parall\`eles, tangente horizontale en $O$, sym\'etrie centrale.
\item \begin{enumerate}
      \item On multiplie num\'erateur et d\'enominateur par $\expo{-x}$. Avec $u(x) = 1 + \expo{-x}$, $u'(x) = -\expo{-x}$ :
            $\frac{4\expo{-x}}{1 + \expo{-x}} = -4\frac{u'(x)}{u(x)}$, une primitive est $x \mapsto -4\ln(1 + \expo{-x})$.
      \item $(\mathscr{C})$ est au-dessus de $(D_1)$ : $\mathscr{A}(\lambda) = \displaystyle\int_0^\lambda \frac{4}{\expo{x} + 1}\,\mathrm{d}x
            = 4\ln 2 - 4\ln(1 + \expo{-\lambda})$ cm$^2$, et $\lim\limits_{\lambda \to +\infty} \mathscr{A}(\lambda) = 4\ln 2 \approx 2{,}77$~cm$^2$.
      \end{enumerate}
\item $f$ est une bijection continue strictement croissante de $\mathbb{R}$ sur $\mathbb{R}$, donc $f^{-1}$ existe sur $\mathbb{R}$.
      Comme $f(0) = 0$ et $f'(0) = 0$, $f^{-1}$ n'est pas d\'erivable en $0$ (tangente verticale en $O$ \`a sa courbe).
\end{enumerate}

%%%%%%%%%%%%%%%%%%%%%%%%%%%%%%%%%%%%%%%%%%%%%%%%%%%%%%%%%%%%%%%%%%%%%%%%%%%%%%%%%%%%%%%%%%%%%%%%%%%%%%%%
\subsection{Corrig\'e du sujet D6}
\noindent \underline{\textsc{Exercice I}}
\begin{enumerate}
\item \begin{enumerate}
      \item $1 + \expo{i\theta} = \expo{i\frac{\theta}{2}}\left(\expo{-i\frac{\theta}{2}} + \expo{i\frac{\theta}{2}}\right) = 2\cos\left(\frac{\theta}{2}\right)\expo{i\frac{\theta}{2}}$.
      \item Pour $\theta \in\ ]-\pi ; \pi[$, $\frac{\theta}{2} \in\ ]-\frac{\pi}{2} ; \frac{\pi}{2}[$ et $\cos\frac{\theta}{2} > 0$ :
            $|z| = 2\cos\frac{\theta}{2}$ et $\arg z = \frac{\theta}{2}$ $[2\pi]$.
      \item $\theta = \frac{2\pi}{3}$ : $1 + \expo{i\frac{2\pi}{3}} = 2\cos\frac{\pi}{3}\,\expo{i\frac{\pi}{3}} = \expo{i\frac{\pi}{3}}$.
            $\theta = \frac{\pi}{2}$ : $1 + i = 2\cos\frac{\pi}{4}\,\expo{i\frac{\pi}{4}} = \sqrt{2}\,\expo{i\frac{\pi}{4}}$.
      \end{enumerate}
\item \begin{enumerate}
      \item $S : z' = az + b$ avec $a = \frac{z_{B'} - z_{A'}}{z_B - z_A} = \frac{-2}{-1 + i} = \frac{-2(-1 - i)}{2} = 1 + i$ et
            $b = z_{A'} - a z_A = 1 + i - (1 + i) = 0$ : $z' = (1 + i)z$.
      \item Centre $O$, rapport $\sqrt{2}$, angle $\frac{\pi}{4}$.
      \end{enumerate}
\item \begin{enumerate}
      \item $z_0 = 1$ et $z_{n+1} = (1 + i)z_n$ : par r\'ecurrence, $z_n = (1 + i)^n = (\sqrt{2})^n\expo{i\frac{n\pi}{4}}$.
      \item $z_{12} = 2^6\expo{3i\pi} = -64$. $M_n \in [O ; -\vec{u})$ si et seulement si $\frac{n\pi}{4} \equiv \pi$ $[2\pi]$,
            c'est-\`a-dire $n = 4 + 8k$, $k \in \mathbb{N}$.
      \end{enumerate}
\end{enumerate}

\noindent \underline{\textsc{Exercice II}} \newline
$n = 6$, $\sum x_i = 57$, $\sum y_i = 131$, $\sum x_i^2 = 689$, $\sum y_i^2 = 2\,933$, $\sum x_i y_i = 1\,141$.
\begin{enumerate}
\item Le nuage descend r\'eguli\`erement ; $G(9{,}5 ; 21{,}83)$ (exactement $\bar{y} = \frac{131}{6}$).
\item $V(x) = \frac{689}{6} - 9{,}5^2 \approx 24{,}583$ ; $V(y) = \frac{2\,933}{6} - \left(\frac{131}{6}\right)^2 \approx 12{,}139$ ;
      $\sigma_{xy} = \frac{1\,141}{6} - 9{,}5 \times \frac{131}{6} = -17{,}25$ ; $r \approx -0{,}999$.
      $r < 0$ : la temp\'erature diminue quand l'altitude augmente ; la corr\'elation est tr\`es forte.
\item $a = \frac{-17{,}25}{24{,}583} \approx -0{,}70$ ; $b = \bar{y} - a\bar{x} \approx 28{,}50$ : $\boxed{y = -0{,}70\,x + 28{,}50}$.
      La temp\'erature baisse d'environ $0{,}7^\circ$C chaque fois que l'on monte de 100~m.
\item 2\,000~m correspond \`a $x = 20$ : $y \approx -14 + 28{,}50 = 14{,}5^\circ$C.
\item $a' = \frac{\sigma_{xy}}{V(y)} \approx -1{,}421$ et $x - 9{,}5 = a'(y - \bar{y})$, soit $x \approx -1{,}421\,y + 40{,}53$.
      Pour $y = 15$ : $x \approx 19{,}2$, c'est-\`a-dire une altitude d'environ 1\,920~m.
\end{enumerate}

\noindent \underline{\textsc{Probl\`eme}}
\begin{enumerate}
\item En $-\infty$, $\expo{x} \to 0$ et $\expo{2x} \to 0$ : $f(x) \to 2$, la droite $(\Delta) : y = 2$ est asymptote en $-\infty$.
      En $+\infty$ : $f(x) = \expo{2x}\left(1 - 3\expo{-x} + 2\expo{-2x}\right) \to +\infty$ et $\frac{f(x)}{x} = \frac{\expo{2x}}{x}(1 - 3\expo{-x} + 2\expo{-2x}) \to +\infty$ :
      branche parabolique de direction $(Oy)$.
\item $f'(x) = 2\expo{2x} - 3\expo{x} = \expo{x}(2\expo{x} - 3)$, du signe de $2\expo{x} - 3$ : $f$ est d\'ecroissante sur
      $\left]-\infty ; \ln\frac{3}{2}\right]$ et croissante sur $\left[\ln\frac{3}{2} ; +\infty\right[$. Minimum :
      $f\left(\ln\frac{3}{2}\right) = \frac{9}{4} - \frac{9}{2} + 2 = -\frac{1}{4}$.
\item $(\expo{x} - 1)(\expo{x} - 2) = \expo{2x} - 3\expo{x} + 2$. $f(x) = 0 \iff x = 0$ ou $x = \ln 2$ ;
      $f(x) < 0$ sur $]0 ; \ln 2[$ et $f(x) > 0$ sur $]-\infty ; 0[ \,\cup\, ]\ln 2 ; +\infty[$.
\item $f(x) - 2 = \expo{x}(\expo{x} - 3)$ : $(\mathscr{C})$ est au-dessous de $(\Delta)$ sur $]-\infty ; \ln 3[$, au-dessus sur
      $]\ln 3 ; +\infty[$ ; ils se coupent en $(\ln 3 ; 2)$.
\item $f(0) = 0$, $f'(0) = -1$ : $(T) : y = -x$.
\item Sur $[0 ; \ln 2]$, $f \leq 0$ : l'aire vaut
      $-\displaystyle\int_0^{\ln 2} f(x)\,\mathrm{d}x = -\left[\frac{\expo{2x}}{2} - 3\expo{x} + 2x\right]_0^{\ln 2}
      = -\left[(2 - 6 + 2\ln 2) - \left(\frac{1}{2} - 3\right)\right] = \frac{3}{2} - 2\ln 2$ unit\'es d'aire,
      soit $4\left(\frac{3}{2} - 2\ln 2\right) = \boxed{(6 - 8\ln 2) \text{ cm}^2} \approx 0{,}45$~cm$^2$.
\item \begin{enumerate}
      \item $g$ est continue et strictement croissante sur $I$ : bijection de $I$ sur $J = \left[-\frac{1}{4} ; +\infty\right[$.
      \item Pour $y \in J$, on r\'esout $g(x) = y$ ; avec $X = \expo{x} \geq \frac{3}{2}$ : $X^2 - 3X + 2 - y = 0$,
            $\Delta = 1 + 4y \geq 0$, et la seule racine $\geq \frac{3}{2}$ est $X = \frac{3 + \sqrt{1 + 4y}}{2}$.
            Donc $\boxed{g^{-1}(y) = \ln\left(\frac{3 + \sqrt{4y + 1}}{2}\right)}$.
      \end{enumerate}
\end{enumerate}
\end{corrige}
